\section{Background}
\label{sec:background}

\subsection{How the detectors track memory}

\paragraph{AddressSanitizer.} ASan~\cite{serebryany2012asan} is compiler instrumentation. One shadow
byte describes eight bytes of application memory and says how many of them are addressable. ASan's
runtime replaces \code{malloc}: every block is surrounded by \emph{redzones} whose shadow is marked
unaddressable, and a freed block is poisoned and kept in a \emph{quarantine} so that its memory is
not handed out again soon. Every load and store is preceded by a shadow check. Large requests are
served by a separate path that maps each block on its own. ASan therefore knows each block's
requested size and whether it is live, but only for blocks it allocated itself. For memory it did
not allocate, a program can describe addressability through \emph{manual poisoning},
\code{ASAN\_POISON\_MEMORY\_REGION} and \code{ASAN\_UNPOISON\_MEMORY\_REGION}~\cite{asanpoison}.
The interface describes bytes, not objects: there is no event for ``this object was freed''.

\paragraph{Valgrind Memcheck.} Memcheck~\cite{nethercote2007valgrind} runs an unmodified binary on a
synthetic CPU and shadows every byte with an addressability bit and definedness bits. Its own
\code{malloc} places a redzone (16 bytes by default) on each side of a block and delays reuse of
freed blocks. A program can tell Memcheck about objects it manages itself through client
requests~\cite{valgrindmanual}: \code{VALGRIND\_CREATE\_MEMPOOL} declares a pool,
\code{VALGRIND\_MEMPOOL\_ALLOC} and \code{VALGRIND\_MEMPOOL\_FREE} declare an object's birth and
death inside it, and \code{VALGRIND\_MAKE\_MEM\_NOACCESS} marks bytes unaddressable. Unlike ASan's
poisoning, the pool requests model ownership, so Memcheck can report a free of something that was
never allocated or was already freed.

\paragraph{HWASan.} Hardware-assisted ASan follows the memory-tagging idea behind Arm
MTE~\cite{serebryany2018mte,hwasandesign}. Memory is divided into 16-byte granules and each granule
carries an 8-bit tag in shadow memory. \code{malloc} gives each new block a random tag, writes it to
the block's granules and returns a pointer that carries the same tag in its top byte; \code{free}
retags the granules. Every access compares the pointer's tag with the granule's tag. An overflow into
a neighbour or a use after free meets a different tag, without redzones or a quarantine. A block
whose size is not a multiple of 16 ends in a \emph{short granule} whose shadow value records how many
bytes are valid. Detection is probabilistic: two tags coincide with probability about 1/256. The
hardware must ignore the tag bits when it dereferences a pointer, which AArch64 does (Top Byte
Ignore) and x86-64 does only with Linear Address Masking (LAM). Allocators can tag their own memory
with \code{\_\_hwasan\_tag\_memory} and \code{\_\_hwasan\_tag\_pointer}; this interface works on whole
granules only.

\paragraph{UBSan and LowFat.} The proposal also named UBSan's \code{bounds} and \code{object-size}
checks and LowFat~\cite{duck2016lowfat}. UBSan inserts local checks that rely on sizes the compiler
can see; it keeps no record of heap blocks at run time. LowFat allocates each size class from its own
large region of virtual memory, so the bounds of an object can be computed from its address alone,
and checks pointer arithmetic against them; it covers spatial errors only. Section~\ref{sec:limits}
explains why neither entered the detection matrix.

\subsection{Custom allocators and why they hide objects}

A CMA takes a block from \code{malloc} (or a static array) and manages its interior itself. The
allocators we study cover the common designs:
\begin{itemize}
\item \textbf{Arena (bump) allocator.} Objects are carved one after another from the block; freeing
  one object does not make its memory reusable, and the whole arena is released at once.
\item \textbf{Size-class free list.} Each size class keeps a list of free slots; \code{free} pushes a
  slot and the next request of that class pops it, so reuse is immediate. The link to the next free
  slot is stored inside the free slot.
\item \textbf{Buddy allocator} (SQLite memsys5~\cite{sqlitemalloc}). A fixed heap is split into
  power-of-two blocks that are split and coalesced on demand.
\item \textbf{Lookaside cache} (SQLite). Each database connection keeps a slab of small fixed-size
  slots for short-lived objects and falls back to the general allocator for larger ones.
\end{itemize}
To a \code{malloc}-granular detector all of these look the same: one live block. Four properties
decide what it can still see. Objects are adjacent, with at most alignment slack between them and no
redzones. A freed object is still inside a live block, so its memory stays addressable. Under
immediate reuse a stale pointer and the pointer of the new owner hold the same address. And when the
whole pool is released, the detector sees one large \code{free}, whatever objects were inside.
Figure~\ref{fig:layout} contrasts the two views for three 24-byte objects.

\begin{figure}[t]
\centering\footnotesize
\begin{verbatim}
 malloc + ASan    [rz][ a: 24 |8][rz] ... [rz][ b: 24 |8][rz] ... [rz][ c: 24 |8][rz]
                   every gap is poisoned; a freed block is poisoned and quarantined

 Arena            [ a: 24 |8][ b: 24 |8][ c: 24 |8][ ...                    pool ... ]
 (one malloc)      a+24 is slack, a+32 is b[0], a+64 is c[0]: all addressable to the tool

 Arena, annotated [ a: 24 |##][ b: 24 |##][ c: 24 |##][ ##################### ]
                   ## poisoned: the slack and the unused pool; b[0] and c[0] stay live
\end{verbatim}
\vspace{-1em}
\caption{Three \code{st\_alloc(24)} objects. ASan's own \code{malloc} separates them with redzones;
an arena packs them into one block that the tool sees as a single live allocation. Annotation
poisons the slack and the rest of the pool, but cannot poison a live neighbour.}
\label{fig:layout}
\end{figure}
